The \dap{} can play audio through Bluetooth headphones and speakers, which can
also control playback with their buttons. While a Bluetooth device is
connected, audio plays through it instead of the \daps{} own outputs.

\begin{description}
  \item[Bluetooth.] Turns Bluetooth \setting{On} or \setting{Off}. When it is
    on, the \dap{} reconnects to the most recently used device, and paired
    devices can connect to the \dap{} by themselves. Turning Bluetooth off
    saves battery.
  \item[Devices.] Lists the paired devices, followed by the devices found by
    searching for new ones, which starts when the list is opened. New devices
    must be in pairing mode to be found. The first line shows the Bluetooth
    status; select it to search again, or to turn Bluetooth on.

    Select a device to connect to it, or to disconnect it if it is connected.
    Press \ActionStdContext{} on a paired device to forget it.
  \item[High Quality Audio.] When \setting{On}, and the device supports it,
    uses the SBC codec's dual channel mode, which roughly doubles the bitrate.
    Should playback stutter, turn this \setting{Off}.
  \item[Sync Volume with Headphones.] When \setting{On}, and the device
    supports it, the device applies the volume, and changing the volume on the
    device changes it on the \dap{} too. When \setting{Off}, the \dap{} applies
    the volume, independently of the device's own volume.
\end{description}

\opt{headphone_detection}{
When a Bluetooth device disconnects, playback pauses as when unplugging
headphones, following the \setting{Pause on Headphone Unplug} setting (see
\reference{ref:PauseOnUnplug}).
}
